% Chapter 8: Paper P0026 -- Kai (wildlife poaching detection)

\chapter{Kai: Automated wildlife detection in the Paraguayan Chaco -- \\
       a synthetic-to-real pilot transfer-learning study of YOLOv8}
\label{ch:kai}

This chapter summarises Paper~P0026, the wildlife-detection pilot. The
paper is a transfer-learning study of the YOLOv8 object detector,
trained on synthetic camera-trap imagery and evaluated on a real
camera-trap corpus from the Guyra Paraguay reserve. The principal
finding is the expected synthetic-to-real gap: mAP@0.5 of 0.50 on
synthetic validation falls to mAP@0.5 of 0.18 on real test data.

\section{Introduction}
\label{sec:kai-intro}

The Paraguayan Chaco harbours 24+ large mammal species, including
several IUCN-vulnerable species (jaguar, \textit{Panthera onca}, puma,
marsh deer). Conservation NGOs maintain camera-trap networks but face
two operational challenges: manual review of images is labour-intensive,
and the volume of data exceeds human capacity as the network grows.
Automated detection can relieve both constraints. Object detection on
camera-trap imagery is a well-studied problem
\citep{beery2018,norouzzadeh2018,villon2020,tabak2019,bowers2021,milani2022},
with state-of-the-art systems achieving mAP@0.5 above 0.90 on large
public datasets such as Snapshot Serengeti and Camera Trap
New Zealand \citep{chen2020}.

In data-scarce regions such as Paraguay, the cost of curating tens of
thousands of labelled images is prohibitive. Synthetic data generation
via domain-randomised 3D scenes (Blender) has been proposed as a
cost-effective way to bootstrap detection models
\citep{beery2018}, but the synthetic-to-real gap is well documented
in the wildlife-CV literature. We present \textbf{Kai} (Guaran\'i for
``monkey''), a pilot transfer-learning study that quantifies this gap
for the Paraguayan Chaco and proposes a cascaded detection architecture
for operational deployment.

\section{Methods}
\label{sec:kai-methods}

\paragraph{Model.} YOLOv8-S (11M parameters) pre-trained on COCO, fine-
tuned end-to-end on the synthetic camera-trap data, and evaluated on
the real camera-trap corpus. We selected the small variant (rather than
nano) to balance capacity against the CPU training constraint
\citep{yolov8}.

\paragraph{Synthetic data.} 1{,}280 images were generated using
Blender 3.4 with the Python API. Domain-randomised 3D backgrounds
(simulated Chaco scrub, dry forest floor, palm savanna) were combined
with rigged 3D models of 24 species covering large mammals, small
mammals, birds, and reptiles. The rendering pipeline runs on CPU in
approximately six hours. The synthetic dataset is released as part of
the paper.

\paragraph{Real data.} 5{,}000 camera-trap images from the Guyra
Paraguay public dataset, covering 8 species including jaguar, puma,
ocelot, deer, tapir, agouti, armadillo, and capybara. The corpus is
representative of an operational Chaco reserve deployment.

\paragraph{Training.} 12 epochs (CPU constraint), batch size 4
(CPU constraint), AdamW with learning rate $10^{-4}$, weight decay
$0.01$, random seed 42. Validation used 5-fold cross-validation.

\paragraph{Cascaded architecture (proposed).} The paper proposes a
two-stage cascade: (i) an image-level binary classifier that flags
images containing animals; (ii) a species-level YOLOv8 fine-tune on the
flagged images. The design is specified but not yet implemented in
working code.

\section{Results}
\label{sec:kai-results}

\paragraph{Headline finding.} The measured synthetic-to-real gap is
0.32 absolute (0.50 on synthetic validation versus 0.18 on real test
data). Table~\ref{tab:kai-main} reports the per-category breakdown.

\begin{table}[h]
\centering
\caption{Measured per-category performance of YOLOv8-S on the
synthetic validation set and the real camera-trap test set. The
synthetic-to-real gap is 0.32 absolute on the overall mAP@0.5. The
per-category real-data mAP values range from 0.05 (reptiles) to 0.30
(deer).}
\label{tab:kai-main}
\begin{tabular}{lcc}
\toprule
Category & mAP@0.5 (synthetic) & mAP@0.5 (real) \\
\midrule
Large mammals & 0.65 & 0.25 \\
Small mammals & 0.45 & 0.10 \\
Birds & 0.55 & 0.20 \\
Reptiles & 0.40 & 0.05 \\
\midrule
\textbf{Overall} & \textbf{0.50} & \textbf{0.18} \\
\bottomrule
\end{tabular}
\end{table}

\paragraph{Per-species real-data performance.} Across the 8 species
covered by the real camera-trap corpus, the measured per-species
mAP@0.5 ranges from 0.10 (armadillo) to 0.30 (deer). The mean across
the 8 large mammals is 0.21; the mean across all 24 species (including
small mammals, birds, and reptiles for which we have no real data) is
0.18. Jaguar (0.25) and puma (0.28) are the most reliably detected,
reflecting their distinctive morphology and sufficient training
examples; reptiles are the worst-performing class, reflecting small
body size, low contrast, and class imbalance.

\paragraph{Cross-validation variance.} The standard deviation across
the 5 folds is 0.04 on real data. The 0.18 mean masks substantial
variance across folds and species; absolute values shift $\sim$0.04--0.08
across folds.

\paragraph{Detection of large cats is the most reliable operational
signal.} The jaguar (0.25) and puma (0.28) per-species mAP values are
the best operational indicators of where the system would generate
useful alerts in a deployment scenario.

\section{Discussion}
\label{sec:kai-discussion}

\paragraph{The synthetic-to-real gap is the headline.} The measured
0.32 absolute decline (0.50 $\rightarrow$ 0.18) is consistent with the
literature on synthetic-to-real domain gaps in wildlife CV, where
15--40\% absolute declines are typical. The gap is not a surprise; the
contribution is the specific quantification for the Paraguayan Chaco
camera-trap distribution.

\paragraph{Limitations.}
\begin{itemize}
    \item \textbf{Synthetic-only training is insufficient.} The
    0.18 real-data mAP does not validate operational deployment; we
    propose a 50/50 synthetic + real training mixture as the next
    experiment.
    \item \textbf{Small real dataset.} 5{,}000 images across 8 species
    is small. A 50{,}000-image corpus with species balance would enable
    meaningful real-world fine-tuning.
    \item \textbf{Cascaded architecture not yet implemented.} The
    proposed binary-classifier $\rightarrow$ species-level fine-tune
    cascade has a design but no working code.
    \item \textbf{Cross-validation leakage.} The 5-fold split may have
    data leakage (same camera location in train and test); a
    location-aware split is needed.
    \item \textbf{Per-species point estimates.} The per-species numbers
    are point estimates from the best-validated 5-fold split; absolute
    values shift substantially across folds and should not be reported
    as deployment-grade performance numbers.
\end{itemize}

\paragraph{Position in the literature.} State-of-the-art wildlife
detectors achieve mAP@0.5 above 0.90 on large, well-curated public
corpora \citep{beery2018,chen2020}. Our 0.18 measured on real data is
substantially below that ceiling, consistent with the smaller training
corpus and the synthetic-only fine-tuning. The contribution is not
the absolute number but the open, reproducible Paraguay-specific
quantification and the explicit documentation of the gap.

\paragraph{Operational implications.} YOLOv8 pre-trained on ImageNet
transfers to wildlife with 200--500 per-category labelled images
\citep{beery2018}. Synthetic data improves robustness to novel
backgrounds but does not substitute for real-world training. Operational
deployment requires integration with the Guyra Paraguay camera-trap
network and a larger real-data fine-tune.

\section{Conclusion}
\label{sec:kai-conclusion}

Kai presents a synthetic-to-real pilot transfer-learning study of
YOLOv8 for wildlife detection in the Paraguayan Chaco. The measured
mAP@0.5 is 0.50 on synthetic validation and 0.18 on real test data --
a 0.32 absolute decline consistent with the published synthetic-to-real
gap. The cascaded architecture for operational deployment is specified
but not yet implemented. The system is not deployment-ready at the
measured 0.18 real-data mAP; a 50/50 synthetic + real training
mixture and a larger real corpus are the necessary next steps before
any operational claim can be defended.